% GENERATED FILE. Edit canonical lessons, then run npm run generate:books.
% canonical-source-hash: fnv1a64:905dcf2000d88bb5
% canonical-lessons: RU-C462-palatka, RU-C462-spalnik, RU-C462-kostyor, RU-C462-razzhigat, RU-C462-tushit

\chapter{A night in a tent}
\label{ch:ru-a-night-in-a-tent}
\hlchaptermodality{462}

\begin{chapteropening}
\textbf{By the end of this chapter:} \emph{I can go camping: a tent, a sleeping bag, a campfire, lighting a fire and putting it out.}

Complete the last lesson of chapter 462: I can go camping: a tent, a sleeping bag, a campfire, lighting a fire and putting it out.
\end{chapteropening}

\section[\texorpdfstring{\ru{палатка} (palátka)}{палатка (palátka)}]{\ru{палатка} (palátka) — a tent}

\label{lesson:RU-C462-palatka}



\begin{quote}
Before the new one: say the Russian for moss, then the Russian for a cone.
\end{quote}

\subsection*{You'll want to know: \ru{палатка}}
\textbf{\ru{палатка}} (\emph{palátka}) — \textquotedblleft{}a tent\textquotedblright{}.

\begin{quote}
\ru{Мы} \ru{поставили} \ru{палатку} \ru{у} \ru{реки}.
\end{quote}

\begin{quote}
\emph{My postávili palátku u rekí.}
\end{quote}

\begin{quote}
We put up the tent by the river.
\end{quote}

Feminine. \emph{\ru{Ставить} \ru{палатку}} — to pitch a tent.

\begin{cousinweb}
\emph{\ru{палатка}} is usually traced to \emph{\ru{палата}} 'a chamber, a hall', which goes back through Greek to Latin \emph{palātium} 'a palace' (Lewis \& Short). EtymDB lists \emph{\ru{палатка}} with the sense 'tent' but no source.
\end{cousinweb}

\begin{grammarlens}[title={What you've built}]
One more word for telling, explaining and describing at B1.
\end{grammarlens}

\subsection*{Your turn}
\begin{itemize}
\raggedright
  \item \emph{Say it:} \emph{palátka}
  \item \emph{Say it:} \emph{palátka}, once more
  \item \emph{Say it:} a sentence of your own with \emph{palátka}
  \item \emph{Recall:} say the Russian for moss, then the Russian for a cone, then say \emph{palátka} again
\end{itemize}

\begin{tcolorbox}[breakable,colback=teal!4,colframe=teal!35!black,title={Before you move on}]
What does \ru{палатка} mean? (\textquotedblleft{}A tent\textquotedblright{}.) Say it once more.
\end{tcolorbox}
\section[\texorpdfstring{\ru{спальник} (spálnik)}{спальник (spálnik)}]{\ru{спальник} (spálnik) — a sleeping bag}

\label{lesson:RU-C462-spalnik}



\begin{quote}
Before the new one: say the Russian for a cone, then the Russian for a tent.
\end{quote}

\subsection*{You'll want to know: \ru{спальник}}
\textbf{\ru{спальник}} (\emph{spálnik}) — \textquotedblleft{}a sleeping bag\textquotedblright{}.

\begin{quote}
\ru{Ночью} \ru{в} \ru{спальнике} \ru{было} \ru{тепло}.
\end{quote}

\begin{quote}
\emph{Nóchyu v spálnike býlo tepló.}
\end{quote}

\begin{quote}
At night it was warm in the sleeping bag.
\end{quote}

Colloquial for \emph{\ru{спальный} \ru{мешок}}.

\begin{cousinweb}
\emph{\ru{спальник}} is a colloquial short form of \emph{\ru{спальный} \ru{мешок}} 'a sleeping bag', from \emph{\ru{спать}} 'to sleep', which is inherited from Proto-Slavic \emph{sŭpati} (EtymDB).
\end{cousinweb}

\begin{grammarlens}[title={What you've built}]
One more word for telling, explaining and describing at B1.
\end{grammarlens}

\subsection*{Your turn}
\begin{itemize}
\raggedright
  \item \emph{Say it:} \emph{spálnik}
  \item \emph{Say it:} \emph{spálnik}, once more
  \item \emph{Say it:} a sentence of your own with \emph{spálnik}
  \item \emph{Recall:} say the Russian for a cone, then the Russian for a tent, then say \emph{spálnik} again
\end{itemize}

\begin{tcolorbox}[breakable,colback=teal!4,colframe=teal!35!black,title={Before you move on}]
What does \ru{спальник} mean? (\textquotedblleft{}A sleeping bag\textquotedblright{}.) Say it once more.
\end{tcolorbox}
\section[\texorpdfstring{\ru{костёр} (kostyór)}{костёр (kostyór)}]{\ru{костёр} (kostyór) — a campfire, a bonfire}

\label{lesson:RU-C462-kostyor}



\begin{quote}
Before the new one: say the Russian for a tent, then the Russian for a sleeping bag.
\end{quote}

\subsection*{You'll want to know: \ru{костёр}}
\textbf{\ru{костёр}} (\emph{kostyór}) — \textquotedblleft{}a campfire, a bonfire\textquotedblright{}.

\begin{quote}
\ru{Вечером} \ru{мы} \ru{пели} \ru{песни} \ru{у} \ru{костра}.
\end{quote}

\begin{quote}
\emph{Vécherom my péli pésni u kostrá.}
\end{quote}

\begin{quote}
In the evening we sang songs by the campfire.
\end{quote}

Masculine; the \ru{ё} drops out: \emph{\ru{костра}}, \emph{\ru{костры}}.

\begin{cousinweb}
The origin of \emph{\ru{костёр}} is not settled; EtymDB lists the word without a source.
\end{cousinweb}

\begin{grammarlens}[title={What you've built}]
One more word for telling, explaining and describing at B1.
\end{grammarlens}

\subsection*{Your turn}
\begin{itemize}
\raggedright
  \item \emph{Say it:} \emph{kostyór}
  \item \emph{Say it:} \emph{kostyór}, once more
  \item \emph{Say it:} a sentence of your own with \emph{kostyór}
  \item \emph{Recall:} say the Russian for a tent, then the Russian for a sleeping bag, then say \emph{kostyór} again
\end{itemize}

\begin{tcolorbox}[breakable,colback=teal!4,colframe=teal!35!black,title={Before you move on}]
What does \ru{костёр} mean? (\textquotedblleft{}A campfire, a bonfire\textquotedblright{}.) Say it once more.
\end{tcolorbox}
\section[\texorpdfstring{\ru{разжигать} (razzhigát)}{разжигать (razzhigát)}]{\ru{разжигать} (razzhigát) — to light (a fire), to kindle}

\label{lesson:RU-C462-razzhigat}



\begin{quote}
Before the new one: say the Russian for a sleeping bag, then the Russian for a campfire.
\end{quote}

\subsection*{You'll want to know: \ru{разжигать}}
\textbf{\ru{разжигать}} (\emph{razzhigát}) — \textquotedblleft{}to light (a fire), to kindle\textquotedblright{}.

\begin{quote}
\ru{Папа} \ru{разжигает} \ru{костёр} \ru{одной} \ru{спичкой}.
\end{quote}

\begin{quote}
\emph{Pápa razzhigáyet kostyór odnóy spíchkoy.}
\end{quote}

\begin{quote}
Dad lights the campfire with one match.
\end{quote}

Imperfective; perfective \emph{\ru{разжечь}} (\emph{\ru{разожгу}, \ru{разожжёшь}}).

\begin{cousinweb}
\emph{\ru{разжигать}} is built from \emph{\ru{раз}-} and \emph{-\ru{жигать}}, from \emph{\ru{жечь}} 'to burn', the same root as \emph{\ru{зажигать}} 'to light'. EtymDB lists \emph{\ru{разжигать}} without a deeper source.
\end{cousinweb}

\begin{grammarlens}[title={What you've built}]
One more word for telling, explaining and describing at B1.
\end{grammarlens}

\subsection*{Your turn}
\begin{itemize}
\raggedright
  \item \emph{Say it:} \emph{razzhigát}
  \item \emph{Say it:} \emph{razzhigát}, once more
  \item \emph{Say it:} a sentence of your own with \emph{razzhigát}
  \item \emph{Recall:} say the Russian for a sleeping bag, then the Russian for a campfire, then say \emph{razzhigát} again
\end{itemize}

\begin{tcolorbox}[breakable,colback=teal!4,colframe=teal!35!black,title={Before you move on}]
What does \ru{разжигать} mean? (\textquotedblleft{}To light (a fire), to kindle\textquotedblright{}.) Say it once more.
\end{tcolorbox}
\section[\texorpdfstring{\ru{тушить} (tushít)}{тушить (tushít)}]{\ru{тушить} (tushít) — to put out (a fire); to stew}

\label{lesson:RU-C462-tushit}



\begin{quote}
Before the new one: say the Russian for a campfire, then the Russian for to light.
\end{quote}

\subsection*{You'll want to know: \ru{тушить}}
\textbf{\ru{тушить}} (\emph{tushít}) — \textquotedblleft{}to put out (a fire); to stew\textquotedblright{}.

\begin{quote}
\ru{Перед} \ru{сном} \ru{мы} \ru{тушим} \ru{костёр} \ru{водой}.
\end{quote}

\begin{quote}
\emph{Péred snom my túshim kostyór vodóy.}
\end{quote}

\begin{quote}
Before going to sleep we put out the campfire with water.
\end{quote}

Imperfective; perfective \emph{\ru{потушить}}. Present \emph{\ru{тушу}, \ru{тушишь}}. In cooking: \emph{\ru{тушить} \ru{мясо}} — to stew meat.

\begin{cousinweb}
\emph{\ru{тушить}} goes with \emph{\ru{тухнуть}} 'to go out (of a fire)'. EtymDB lists \emph{\ru{тушить}} without a source, so the deeper origin is left open.
\end{cousinweb}

\begin{grammarlens}[title={What you've built}]
One more word for telling, explaining and describing at B1.
\end{grammarlens}

\subsection*{Your turn}
\begin{itemize}
\raggedright
  \item \emph{Say it:} \emph{tushít}
  \item \emph{Say it:} \emph{tushít}, once more
  \item \emph{Say it:} a sentence of your own with \emph{tushít}
  \item \emph{Recall:} say the Russian for a campfire, then the Russian for to light, then say \emph{tushít} again
\end{itemize}

\begin{tcolorbox}[breakable,colback=teal!4,colframe=teal!35!black,title={Before you move on}]
What does \ru{тушить} mean? (\textquotedblleft{}To put out (a fire); to stew\textquotedblright{}.) Say it once more.
\end{tcolorbox}
